% ch3 D.4 反铁磁自旋波与“对称破缺” (原书页 175–182 / PDF p190–p197)

% 衔接说明：文件最顶部（%--D3-TAIL-- 标记之下）为前一小节 D.3 的收尾文字
%（原书 p.175 顶部、小节 D.4 标题之前的一段，续接 p.174 末句），装配时应
% 移至承接原书 p.174（PDF p189）的 D.3 文件末尾拼接。

%--D3-TAIL--
%JOIN%
1 的反向磁畴出现在样品中的任何地方，四周由厚度 $\sim 1$ 的“Bloch 墙”所包围，其能量为（墙中原子数）$\times\,(\nabla M)^2 \sim 1^2(1/1^2)M^2 =$ 有限值。于是，在样品的所有区域内、在任意尺度上、并且以实质上无限多种可能的组态，这类磁畴都具有一份有限的热激发概率。这样，长程序便被破坏了。一般说来，只要变化缓慢，Landau-Lifshitz 方程使我们能够处理比自旋波更复杂、更一般的运动与组态。

\subsection{反铁磁自旋波与“对称破缺”}

最后，我要对反铁磁自旋波（antiferromagnetic spin waves）这一课题 (101) 作一简短的介绍。它们之所以重要，是因为在零点运动、对称性与元激发谱之间的种种关系中，它们大概是我们能归并在“对称破缺理论”（theory of broken symmetry）名下的各种情形中最简单、也理解得最好的一例；而这些关系在几乎所有形式的量子凝聚——固体与声子、自旋波、超导电性、液氦以及等离体——的理论中都起着至关重要的作用。

一般说来，凝聚现象的特征是：在我们称为“凝聚”系统的状态中出现了一个原先系统中不存在的新参量。铁磁体在 Curie 点以上、没有外场时没有磁化；在 Curie 点以下则有。在反铁磁体中，新参量便是反铁磁序。例如，我们可以取两个子晶格磁化之差作为参量——它常被称为“序参量”（order parameter）。合金中的有序-无序有一个明显的新参量，铁电性也是如此。同样明显的是固体的长程序，它在液体中并不存在。液氦 II 有一个参量，如今普遍理解为处于零动量态中的原子数。液-气系统则更难以捉摸一些，但利用所谓“格气”（lattice gas）方法，我们可以把它类比于一个经典的有序-无序系统。更加难以捉摸、然而却完全类似的，是超导体的序参量，它可以定义为与能隙相联系的某个非对角的“Green 函数”。

所有这些序参量都有一个共同的特点：当序参量出现时，系统比起未凝聚相来具有较低的对称性。这一点在多数情形下又是十分明显的，尽管对两种超流体而言，这种较低的对称性是相当特殊而且有争议的一种，即通常意义下规范不变性的缺失（也许更为普遍接受的观点是零动量起着特殊作用，即缺乏 Newton 不变性）。在多数其他情形中，对称性的改变属于那种对一般理论物理学家的敏感神经刺激较小的一类，因为他更熟悉那些破坏对称性的微扰的存在。这样，理想的铁磁体与反铁磁体都已抛弃了自旋空间中交换力的转动对称性——这种对称性反正会被磁性微扰所破坏。铁磁态与反铁磁态不具有时间反演对称这一事实（对铁磁体而言这是一件非常实际的事，因为它使铁氧体回转器得以存在）本应几乎同超导体中的规范（不变性）一样使理论家感到不安，因为时间反演通常被认为是一种严格成立的对称性。固体已经抛弃了完全的转动对称性与平移对称性，这同样不会扰乱我们的直觉，但这仅仅是因为我们天天都在同固态的夹具和匣子打交道，而它们恰好具有同样奇特的对称性缺失。尽管如此，在所有这些情形中，我们通常都取一个\emph{具有}适当对称性的哈密顿量表达式，然后我们从它导出——或者假定我们能够导出——系统的一个态，该态由这个哈密顿量决定，但却\emph{不具有}这种对称性。这样看来，据我所见，这正是一切凝聚现象的一个普遍特征。

在所有这些凝聚相之中，铁磁性是独一无二的，或者差不多是独一无二的：在这种情形中，序参量至少近似地是一个\emph{运动常量}（constant of the motion）。这一事实使我们能够作极大的简化。首先，在纯铁磁体的理想化 Heisenberg 模型中，我们能够写下系统的一个精确基态；在更复杂的模型中，例如 Heisenberg 亚铁磁体，我们至少能够用一个好量子数——总自旋，和/或其总的 $z$ 分量——来表征基态。这样一个好量子数的存在，意味着关于基态对称性的一个严格论断。于是我们可以这样表征所有激发态：其中至少有一个被反转的自旋。这种情况就像 $N+1$ 体问题中的情形一样：带有元激发的态确实能够同基态清楚地分辨开来。而且，磁化反转的数量也同粒子数的改变一样，是每个元激发的一个固定常量。我们在前面已经指出，哈密顿量中在铁磁态被“破坏”的那个对称性的存在，当然意味着实际上存在着一整族基态——在此时此地为 $N+1$ 个——它们通过整个自旋系统的转动而相互变换。由于自旋是运动常量，这些不同的态都是哈密顿量的精确本征态，而且正如我们已看到的，这意味着 $\omega(k{=}0)$ 恒等于零。

在几乎所有其他有趣的凝聚情形中，序参量都不是运动常量。让我们以最简单的情形——反铁磁性——为例。一种特别简单的理想反铁磁体的哈密顿量是

\begin{equation*}
\mathcal{H} \;=\; J\sum_{jl} \mathbf{S}_j\cdot\mathbf{S}_l
\end{equation*}

其中我们不妨作一些非实质性的简化：设各个 $\mathbf{S}$ 位于一个简单立方晶格上，而 $J$ 只在最近邻之间起作用，于是 $\mathbf{S}_j$ 总可以取在面心立方子晶格 A 上，而 $\mathbf{S}_l$ 在 B 上（图 57）。
%FIG{57}: 图 57　（原书图注仅作 “Figure 57”）简单立方晶格的 3×3 格点阵列：中心格点属 A 子晶格，其自旋以箭头标为 $\mathbf{S}_j$；其上、下、左、右四个近邻格点属 B 子晶格，其一以箭头标出自旋 $\mathbf{S}_l$；四角格点亦属 A 子晶格

A 子晶格自旋的运动方程是

\begin{align*}
\bigl[\mathcal{H},\ \mathbf{S}_A\bigr] \;&=\; J\sum_j \mathbf{S}_j \times \sum_{\text{neighbors to j}} \mathbf{S}_l \\
&=\; J\sum_{\text{nearest-neighbor pairs}} \mathbf{S}_j \times \mathbf{S}_l \;=\; -\bigl[\mathcal{H},\ \mathbf{S}_B\bigr]
\end{align*}

其中 $\mathbf{S}_B$ 是另一子晶格的总自旋。于是，尽管总自旋 $\mathbf{S}_A + \mathbf{S}_B$ 按我们的一般定理是运动常量，但自旋之差 $\mathbf{S}_A - \mathbf{S}_B$ 的任何一个分量却都不是：

\begin{equation*}
\bigl[\mathcal{H},\ \mathbf{S}_A - \mathbf{S}_B\bigr] \;=\; 2J\sum_{\text{pairs}} \mathbf{S}_j \times \mathbf{S}_l
\end{equation*}

右端这个量，在系统邻近反铁磁态的任何可能态中都断然不为零；在铁磁态中它则可以为零。例如，它的（−）分量为

\begin{equation*}
\bigl[\mathcal{H},\ \mathbf{S}_A^- - \mathbf{S}_B^-\bigr] \;=\; J\sum_{j,l} \left(S_j^- S_{lz} - S_l^- S_{jz}\right)
\end{equation*}

在 $S_{jz} = S$、$S_{lz} = -S$ 的态（它决不是一个本征态）中，它就是

\begin{equation*}
\frac{JS}{\hbar}\sum_{jl}\left(S_j^- + S_l^-\right)
\end{equation*}

其矩阵元在该态中并不为零，尽管当然不存在\emph{对角}元。一般说来，我不作冗长的论证，只向读者保证：在反铁磁情形，$\mathbf{S}_A - \mathbf{S}_B$ 不可能是运动常量。

另一方面，毫无疑问确实发生的是：这个“序参量”变量的进动速率可以极其缓慢——在系统尺寸趋于无限的极限下，缓慢到可以忽略。

让我们取严格的运动方程

\begin{equation*}
i\hbar\,\frac{d}{dt}\left(\mathbf{S}_A - \mathbf{S}_B\right) \;=\; 2J\sum_{\text{pairs}}\left[\mathbf{S}_j \times \mathbf{S}_l\right]
\end{equation*}

两个子晶格的平均自旋 $\mathbf{S}_A$ 与 $\mathbf{S}_B$ 是非常多个原子的自旋之和，因而颇像经典变量，诸如宏观物体的位置与动量；它们的不确定度只需具有相对量级 $N^{-1/2}$。如果我们假定 $\mathbf{S}_A$ 与 $\mathbf{S}_B$ 确实具有某种宏观平均值，那么按对称性，量 $\mathbf{S}_j$ 与 $\mathbf{S}_l$ 的平均值——虽然不是它们的非对角矩阵元，即量子涨落——看来应当由下式给出

\begin{equation*}
\langle \mathbf{S}_j\rangle \;=\; \frac{\langle \mathbf{S}_A\rangle}{N}\,,\qquad \langle \mathbf{S}_l\rangle \;=\; \frac{\langle \mathbf{S}_B\rangle}{N}
\end{equation*}

于是，代入各变量的\emph{平均值}（average values），我们可以写出一个 c 数（c-number）运动方程

\begin{equation*}
i\hbar\,\frac{d}{dt}\langle \mathbf{S}_A - \mathbf{S}_B\rangle \;=\; \frac{2JZ}{N}\,\langle \mathbf{S}_A\rangle \times \langle \mathbf{S}_B\rangle
\end{equation*}

取 $\langle \mathbf{S}_A\rangle = -\langle \mathbf{S}_B\rangle$——这显然是平均能量最低的态——右端便变为零，而且至少在 $N$ 的量级上，序参量 $\langle \mathbf{S}_A\rangle - \langle \mathbf{S}_B\rangle$ 成为常量。当然，我们在这个方程右端略去了量子涨落，但我们可以假设，相对于任何 $\sim N$ 个元素之和的平均值，涨落具有 $1/\sqrt{N}$ 的量级。这完全不是一个严格的证明，但足以指示事情的走向。

另一方面，由于 $\mathbf{S}_A - \mathbf{S}_B$ 并非真正的运动常量，它的值必定要涨落。事实上，涨落必定具有 $N$ 的量级。Hulthen (103) 曾证明，对任何具有 $\mathcal{H} = J\sum_{\text{n.n.}}\mathbf{S}_j\cdot\mathbf{S}_l$（若 $J>0$）的有限系统，其真实的精确基态必有 $S_{\text{tot}} = \mathbf{S}_A + \mathbf{S}_B = 0$。这在物理上其实是显而易见的。但在 $S_{\text{tot}} = 0$ 的态中，$\langle \mathbf{S}_A\rangle = \langle \mathbf{S}_B\rangle = 0$，因为这样的态是完全各向同性的。（这是一个严格的群论结果。）

这看起来是一个颇为令人困惑的悖论。观测以及一种推理方式告诉我们 $\langle \mathbf{S}_A\rangle = -\langle \mathbf{S}_B\rangle \sim N$，而且这种情形不随时间发生宏观变化；另一种基于对称性的推理则说 $\langle \mathbf{S}_A\rangle = \langle \mathbf{S}_B\rangle = 0$。答案在于：后者在\emph{精确基态}（exact ground state）中是正确的，但这个精确的态无关紧要，因为许多别的态与它相距无穷小（能量上约 $1/N$）。因此，如果我们取 $\langle \mathbf{S}_A\rangle = -\langle \mathbf{S}_B\rangle \sim N$ 的态作为我们的态，它将维持非常长的时间，但由于它实际上是一个由许多精确态组成的波包，最终将会解体。

这与一个宏观固体在自由空间中的行为之间，存在着非常贴切而且精确的类似。这样的物体有一个质心坐标 $Q$、一个相应的动量 $P$，以及一个哈密顿量 $P^2/2MN$，$N$ 为粒子数。$P=0$ 将是精确的基态，此时 $Q$ 将剧烈地涨落：$Q$ 不可能是运动常量。事实上，当 $P\to 0$ 时 $\langle Q^2\rangle\to\infty$。但是各能级彼此靠得极近，因此，如果我们想把整个物体定域在一段非常小的距离之内——比如一个晶格常数的量级，$\Delta Q\sim 1$——相应的 $P$ 的增加为 $\sim 1$，$E$ 的增加为 $\sim 1/N$。于是这些能级是如此接近简并，以至于为了实用的目的——例如为了研究固体的内部结构——我们可以把 $\langle Q\rangle$ 固定下来，然后在 $\langle Q\rangle$ 从其预设值明显涨落之前去研究内部的自由度。或许一个更直接的例子，是宏观刚性旋转体的取向 $\theta$ 与角动量 $I\omega$。

但是我们必须认识到两点：其一，使 $Q$ 或 $\mathbf{S}_A - \mathbf{S}_B$ 的一切取值彼此等价的那个对称性的存在，使这个变量的宏观运动频率为零。这正好就是我们在铁磁性中关于序参量 $\mathbf{S}$ 所得到的结果，尽管在我们的情形中序参量并不是运动常量。其二，这个参量的\emph{涨落}必须为无穷大，这一事实带来了另一个结果。在短程力的情形，宏观能量只能依赖于序参量的\emph{梯度}：

\begin{equation*}
U \;\propto\; \bigl|\nabla\left(\mathbf{S}_A - \mathbf{S}_B\right)\bigr|^2 \qquad\text{或}\qquad \bigl|\nabla Q\bigr|^2
\end{equation*}

但这将意味着

\begin{equation*}
\Bigl\langle\bigl|\nabla(\mathbf{S}_A - \mathbf{S}_B)\bigr|^2\Bigr\rangle \;=\; k^2\langle S_k^2\rangle \;\propto\; \langle V\rangle \;=\; \frac{\hbar\omega_k}{2}
\end{equation*}

于是

\begin{equation*}
\langle S_k^2\rangle \;=\; \frac{\omega_k}{k^2} \;\to\; \infty \qquad\text{当 } k\to 0
\end{equation*}

因此我们预期 $\omega_k\to 0$，但 $\omega_k/k^2\to\infty$。实际上，在最广为人知的那些情形（反铁磁性、固体与液体中的声子）中，$\omega_k\propto k$。核物质的表面振荡 (104) 则有 $\omega_k\propto k^{3/2}$。

当存在低频模式时，与这类凝聚总是相伴随的最后一个效应，是一种长程力。长程力同零质量（即当 $k\to 0$ 时 $\omega\to 0$）粒子之间不可避免的关联，我在这里没有时间详细讨论，只打算提及两个最著名的例子：一是经由铁磁或反铁磁矩的极化而起作用的 Suhl-deGennes 长程核自旋-自旋相互作用 (105)；其二自然是人所共知的、围绕点源的弹性形变的长程性质，它使晶体中空位等的弹性相互作用仅按 $1/r$ 衰减。

现在，我简略地给出反铁磁自旋波的数学处理。附带说一句，这里关于自旋波与磁性的许多材料，都载于我发表在《Seitzschrift》第 14 卷上的文章中。反铁磁自旋波在 Nagamiya、Kubo 与 Yosida (101) 处已有相当好的覆盖。Nagamiya 等人提出了一个非常有用的技巧：即把 B 子晶格上自旋的坐标轴相对于 A 子晶格加以反转：对 $S_l$ 而言，$y\to -y$，$z\to -z$，$x\to x$，于是

\begin{align*}
S_{lz} \;&\to\; -S_{lz}\\
S_{lx} \;&\to\; S_{lx}\\
S_{ly} \;&\to\; -S_{ly}
\end{align*}

因此，

\begin{equation*}
\frac{S_{lx}+iS_{ly}}{\sqrt{2}} \;=\; S_l^+ \longrightarrow S_l^-\,,\qquad S_l^- \longrightarrow S_l^+
\end{equation*}

把 $z$ 与 $y$ 加以反转而不反转 $x$，理由是这样一来这个变换就是一个真转动（proper rotation），从而对新坐标轴对易规则依然成立。用新坐标轴表出，

\begin{equation*}
\mathcal{H} \;=\; -J\sum_{\text{jl neighbors}} \left(S_{jz}S_{lz} - S_{jx}S_{lx} + S_{jy}S_{ly}\right)
\end{equation*}

利用近似展开式

\begin{equation*}
S_{jz} \;=\; S - \frac{S_j^- S_j^+}{2S}
\end{equation*}

——这一展开式我们在铁磁情形中已经讨论过——便得到

\begin{equation*}
\mathcal{H} \;\cong\; -NZJS^2 + \frac{J}{2}\sum_{jl}\left(S_j^- S_j^+ + S_l^- S_l^+ + S_j^+ S_l^+ + S_j^- S_l^-\right)
\end{equation*}

当我们 Fourier 变换到自旋波算符（它们现在是近似玻色子）时，

\begin{equation*}
\beta_\lambda \;=\; \frac{1}{\sqrt{NS}}\sum_j \mathrm{e}^{i\lambda\cdot j}\, S_j^+
\end{equation*}

它就变为

\begin{equation*}
\mathcal{H} \;=\; -NZJS^2 + \frac{JS}{2}\sum_\lambda \Bigl\{ Z\left(\beta_\lambda^\dagger\beta_\lambda + \beta_{-\lambda}^\dagger\beta_{-\lambda}\right) + \sum_l \mathrm{e}^{i\lambda(l-j)}\left(\beta_\lambda^\dagger\beta_{-\lambda}^\dagger + \beta_\lambda\beta_{-\lambda}\right)\Bigr\}
\end{equation*}

这种形式的哈密顿量我们从激子问题起就已熟悉；如果读者回头查看那个情形中“反向图”（backwards diagrams）的计算，就会发现自旋波频率是两个系数的平方之\emph{差}（difference）的平方根：

\begin{align*}
\omega_\lambda \;&=\; \frac{1}{2}JS\sqrt{Z^2 - \left(\sum_l \mathrm{e}^{i\lambda(l-j)}\right)^2}\\
&\cong\; \frac{1}{2}JS\sqrt{\frac{Z}{3}}\ \lambda
\end{align*}

把指数函数展开后，上式即告成立。这样，在这一情形中我们确实发现预期 $\omega_\lambda\sim\lambda$ 得到了验证。我们还发现，\emph{零点运动}（zero-point motion）当 $\lambda\to 0$ 时发散；普通振幅 $\beta_\lambda$ 可以用本征解 $B_\lambda$ 写成

\begin{equation*}
\beta_\lambda \;=\; u_\lambda B_\lambda - v_\lambda B_{-\lambda}^\dagger\,,\qquad \text{etc.}
\end{equation*}

其中

\begin{equation*}
u_\lambda \;=\; \cosh\frac{\theta_\lambda}{2}\,,\qquad v_\lambda \;=\; \sinh\frac{\theta_\lambda}{2}
\end{equation*}

而

\begin{equation*}
\tanh\theta_\lambda \;=\; \frac{1}{Z}\sum_{\text{l neighbor to j}} \mathrm{e}^{i\lambda(l-j)} \;\to\; 1 \qquad\text{当 } \lambda\to 0
\end{equation*}

于是 $u$ 与 $v$ 都发散。在二维和三维两种情形，$\sum_\lambda\langle\beta_\lambda^2\rangle$ 中的发散都是可积的，因此人们发现，除了在名义能量 $-NZJS^2$ 之上获得一定的能量收益，以及整个系统绕完美反铁磁态振荡的一份额外的零点能量之外，并没有别的后果。我们必须把这类系统中出现的发散当作一个时间尺度问题来思考。实际的 $\lambda = 0$ 模的确具有发散的振幅，这对于系统能够取严格各向同性的基态来说是必不可少的；但随着 $N\to\infty$，这一零点运动的频率实际上可以忽略，也就是说，借用 Bogolyubov (105) 的说法，各个可能的态是“准简并”（quasi-degenerate）的。因此，一个合适的近似做法是：取 $\lambda = 0$ 模的围绕某一预先固定取向的一包运动解，然后在我们所构成的这一波包的运动破坏单向态的相干性之前的那段时间间隔内，研究所有其余自旋进动简正模的频率以及（收敛的）零点运动。

当然，实际上各向异性总是存在的，它导致一个场项，把 $\beta_\lambda^\dagger\beta_\lambda$ 项的大小提升到（Z + 各向异性），于是频率变为

\begin{equation*}
\omega \;\sim\; \sqrt{\left(JZ + H_{\text{an}}\right)^2 - \left(JZ\right)^2} \;\sim\; \sqrt{H_{\text{an}}H_{\text{ex}}} \;=\; H_{\text{ae}}
\end{equation*}

这就是著名的反铁磁共振频率的“能隙”（energy gap）表达式。它的实际含义并不是说在这种情形下单向态就真的是基态，而是说系统必须通过隧道贯穿、而不是转动的方式，才能到达时间反演且自旋反转的态——这是一个更为缓慢的过程，当然完全可以忽略。在这个系统中，准简并依然存在，但只存在于有限数目的态之间。

所有这些同样的讨论也直接适用于固体中的声子系统。由于转动与平移自由度必须具有发散的零点振幅，频率正比于 $|k|$，而较低模式的零点振幅发散。类似的现象也常见于 $\mathrm{He}_{\mathrm{II}}$（液氦 II）以及一个假想的中性超导体。

值得注意的是：一般而言，一个\emph{连续的}（continuous）破缺对称——平移、规范、转动——导致 $\omega\to 0$，而且通常当 $k\to 0$ 时零点振幅 $\to\infty$；而一个分立的破缺对称——例如磁性情形中的时间反演，或有序-无序情形中子晶格的互换——则不必对凝聚系统的集体模式有任何特殊的后果，除非碰巧至少还存在一个近似的连续对称——磁性情形正是如此，铁电情形碰巧也是如此。

“对称破缺”这整个课题，作为理论物理学的一个形式分支还相当新颖；尽管如我已提到的，其中许多观念早在几十年前就已在具体情形中得到理解。然而我相信，从固态（solid-state）观点出发对它的一般性讨论，在其他地方尚不存在。当进而考虑运动的集体模式同长波电磁效应——Coulomb 力与长程偶极力——之间的相互作用时，还会接踵出现更进一步的复杂情况。例子有自旋波上的静磁效应 (106)，以及超导电性中的等离体效应 (107)。
